\subsection{Characters of a commutative unital algebra}

DEFINITION 4. --- Let A be a commutative unital algebra over K. A unital character is a unital morphism from A to K.

When no confusion can arise, one simply says character instead of unital character. The set of unital characters of A is denoted X(A). If A is the zero algebra, then X(A) is empty.

Let A and B be commutative unital algebras over K and h a unital morphism from A to B. The map \(\chi \mapsto \chi \circ h\) of \(\mathsf{X}(\mathsf{B})\) into \(\mathsf{X}(\mathsf{A})\) is denoted \(\mathsf{X}(h)\) . If k is a morphism from B to a commutative unital algebra, one has \(\mathsf{X}(k \circ h) = \mathsf{X}(h) \circ \mathsf{X}(k)\) . The map \(\mathsf{X}(\mathrm{Id}_{\mathrm{A}})\) is the identity map of \(\mathsf{X}(\mathrm{A})\) .

If h is surjective, \(\mathsf{X}(h)\) is a bijection from \(\mathsf{X}(\mathsf{B})\) onto the set of characters of A that vanish on the kernel of h.

Let \((\mathrm{A}_{1},e_{1}),\ldots,(\mathrm{A}_{n},e_{n})\) be commutative unital algebras over K and let A be the unital algebra \(A_{1}\times\cdots\times A_{n}\) , with identity element \((e_{1},\ldots,e_{n})\) . For each i, identify \(A_{i}\) with an ideal of A and let \(\pi_{i}\) be the canonical map from A onto \(A_{i}\) . Then \(\mathsf{X}(\pi_{i})\) is a bijection from \(\mathsf{X}(\mathrm{A}_{i})\) onto the set \(X_{i}\) of characters of A vanishing on \(\prod_{j\neq i}A_{j}\) . The sets \(X_{i}\) are pairwise disjoint. On the other hand, let \(\chi\in\mathsf{X}(\mathrm{A})\) . Since \(1=\sum\chi(e_{i})\) , there exists i such that \(\chi(e_{i})\neq0\) . For every \(j\neq i\) and every \(y\in A_{j}\) , one has \(\chi(e_{i})\chi(y)=\chi(e_{i}y)=\chi(0)=0\) , hence \(\chi(\mathrm{A}_{j})=0\). Thus, \(\chi\) vanishes on \(\prod_{j\neq i}A_{j}\) , so that \(\mathsf{X}(\mathrm{A})\) is the union of the \(X_{i}\) .

Let B be the unital algebra \(A_{1} \otimes \cdots \otimes A_{n}\) . Denote by \(h_{i}\) the canonical morphism \(A_{i} \to B\) . Then

\[
\chi \mapsto (\chi \circ h _ {1}, \dots , \chi \circ h _ {n})
\]

is a map from \(\mathsf{X}(\mathsf{B})\) into \(\mathsf{X}(\mathsf{A}_1)\times \dots \times \mathsf{X}(\mathsf{A}_n)\) , and

\[
\left(\chi_ {1}, \dots , \chi_ {n}\right) \mapsto \chi_ {1} \otimes \dots \otimes \chi_ {n}
\]

is a map from \(\mathsf{X}(\mathsf{A}_{1}) \times \cdots \times \mathsf{X}(\mathsf{A}_{n})\) into \(\mathsf{X}(\mathsf{B})\) . One verifies that these maps are inverse bijections of one another, by which one identifies \(\mathsf{X}(\mathsf{B})\) with \(\mathsf{X}(\mathsf{A}_{1}) \times \cdots \times \mathsf{X}(\mathsf{A}_{n})\) .

Let A be a commutative unital algebra over K. Let Y be the set of ideals of codimension 1 of A. For every \(\chi \in \mathsf{X}(\mathsf{A})\) , one has \(\operatorname{Ker}(\chi) \in \mathsf{Y}\) . The map \(\chi \mapsto \operatorname{Ker}(\chi)\) is a bijection from \(\mathsf{X}(\mathsf{A})\) onto Y. Indeed, if \(I \in Y\) , there exists a unique isomorphism of the unital K-algebra A/I onto K and the composite morphism

\[
\mathrm{A} \longrightarrow \mathrm{A} / \mathrm{I} \longrightarrow \mathrm{K}
\]

is the unique character of A with kernel I.

DEFINITION 5. --- Let A be a commutative unital algebra over K. For every \(x \in \mathrm{A}\) , we denote by \(\mathcal{G}_{\mathrm{A}}(x)\) , or simply \(\mathcal{G}(x)\) , the map \(\chi \mapsto \chi(x)\) from \(\mathrm{X}(\mathrm{A})\) to K. It is called the Gelfand transform of \(x\) .

The map \(\mathcal{G}\) is a unital morphism from A into the unital algebra \(\mathrm{K}^{\mathsf{X}(A)}\) of maps from \(\mathsf{X}(A)\) to K. It is called the Gelfand transform of A.

If \(x \in \mathrm{A}\) , the image of the Gelfand transform \(\mathcal{G}_{\mathrm{A}}(x)\) of \(x\) is contained in \(\mathrm{Sp}_{\mathrm{A}}(x)\) . Indeed, let \(\chi \in \mathsf{X}(\mathrm{A})\) ; since \(\chi(x - \chi(x)e) = 0\) , the element \(x - \chi(x)e\) is not invertible.

Let B be a unital commutative algebra over K and h a unital morphism from A to B; then \(\mathsf{X}(h)\colon\mathsf{X}(\mathsf{B})\to\mathsf{X}(\mathsf{A})\) defines a unital morphism \(h_{*}\colon\mathsf{K}^{\mathsf{X}(\mathsf{A})}\to\mathsf{K}^{\mathsf{X}(\mathsf{B})}\) , and the diagram:

\[
\begin{array}{l} \text {A} \xrightarrow {\mathcal {G} _ {\text {A}}} \text {K} ^ {\mathsf {X} (\text {A})} \\ \Big \downarrow h \qquad \qquad \qquad \Big \downarrow h _ {*} \\ \text {B} \xrightarrow {\mathcal {G} _ {\text {B}}} \text {K} ^ {\mathsf {X} (\text {B})} \end{array}\tag{3}
\]

is commutative. Indeed, for every \(x \in A\) and every \(\chi \in \mathsf{X}(\mathsf{B})\) , we have:

\[
\begin{array}{r l} & {\mathcal {G} _ {\mathrm{B}} (h (x)) (\chi) = \chi (h (x)) = (\chi \circ h) (x)} \\ & {\qquad = (\mathsf {X} (h) (\chi)) (x)} \\ & {\qquad = \mathcal {G} _ {\mathrm{A}} (x) (\mathsf {X} (h) (\chi))} \\ & {\qquad = h _ {*} (\mathcal {G} _ {\mathrm{A}} (x)) (\chi).} \end{array}\tag{4}
\]

Suppose now that K is a topological field. We then equip \(\mathsf{X}(\mathsf{A})\) with the topology of pointwise convergence on A (cf.~EVT, III, p.~14, example 1), and the topological space \(\mathsf{X}(\mathsf{A})\) is called the character space of A. The topology of \(\mathsf{X}(\mathsf{A})\) is therefore the coarsest for which the functions \(\mathcal{G}_{\mathrm{A}}(x)\) for \(x \in A\) are continuous, and the map \(\chi \mapsto (\chi(a))_{a \in \mathrm{A}}\) identifies the space \(\mathsf{X}(\mathsf{A})\) with a subset of \(K^{A}\) .

When K = R or C, this topology is the topology induced on \(\mathsf{X}(\mathsf{A}) \subset \mathsf{A}^{*}\) by the weak topology \(\sigma(\mathsf{A}^{*}, \mathsf{A})\) on \(A^{*}\) (EVT, II, p.~45, def. 2); as such, we will also say that it is the weak topology on \(\mathsf{X}(\mathsf{A})\) .

If \(h\) is a unital morphism from A to B, the map \(\mathsf{X}(h)\colon \mathsf{X}(\mathsf{B})\to \mathsf{X}(\mathsf{A})\) is continuous. If \(h\) is surjective, the image of \(\mathsf{X}(h)\), namely, the set of characters of A vanishing on the kernel of \(h\), is closed in \(\mathsf{X}(\mathsf{A})\) ; on the other hand, the topology on \(\mathsf{X}(h)(\mathsf{X}(\mathsf{B}))\) deduced from that of \(\mathsf{X}(\mathsf{B})\) by the bijection \(\mathsf{X}(h)\) is the topology of pointwise convergence in A, that is, the topology induced by that of \(\mathsf{X}(\mathsf{A})\) . In other words, \(\mathsf{X}(h)\) is a homeomorphism from \(\mathsf{X}(\mathsf{B})\) onto a closed subset of \(\mathsf{X}(\mathsf{A})\) .

If \(A_{1},\ldots,A_{n}\) are unital commutative algebras over K, the space \(\mathsf{X}(\mathsf{A}_{1}\times\cdots\times\mathsf{A}_{n})\) is thus identified with the topological sum space of \(\mathsf{X}(\mathsf{A}_{1}),\ldots,\mathsf{X}(\mathsf{A}_{n})\) . Likewise, \(\mathsf{X}(\mathsf{A}_{1}\otimes\cdots\otimes\mathsf{A}_{n})\) is identified with the product topological space \(\mathsf{X}(\mathsf{A}_{1})\times\cdots\times\mathsf{X}(\mathsf{A}_{n})\) .

